\documentclass[a4paper]{article}
\usepackage{graphicx, lmodern, geometry, xcolor, enumitem, booktabs, todonotes}
\usepackage[utf8]{inputenc}
\usepackage[T1]{fontenc}
\usepackage[british]{babel}
\usepackage{mathtools, amsopn, amssymb, amsfonts, amsthm}
\usepackage{hyperref}

\geometry{top=3cm, bottom=3cm, headheight=1cm, headsep=2cm, width=14cm, bindingoffset=0mm, marginpar=25mm, marginparsep=5mm}%,showframe}
\setlength{\parindent}{4mm}

\setlist{noitemsep} % no extra space between items
\setlist[enumerate,1]{label = \roman*),ref = \roman*}

\hypersetup{
    breaklinks = {true}, % breaks links over lines
    pdftitle = {BSD Video script},
    pdfauthor = {Christian Wuthrich}
}

\newcommand{\QQ}{\mathbb{Q}}


\begin{document}
  \title{Script for the BSD Video}
  \author{cw}
  \maketitle

\section{Elliptic Curves}
\subsection{I will explain BSD}

\begin{enumerate}
  \item st+te centred on background
  \item st: I have heard about the Riemann hypothesis (bubbles with $\zeta(s)$)
  You are a number theorist. Are you working on this?
  \item te: Not directly. My research concerns another important conjecture,
  also worth a million dollars (bubble $\zeta(s)$ is kicked out for an elliptic curve)
  \item st: What is it?
  \item te: The Birch and Swinnerton-Dyer conjecture.
  \item st: Can you explain?
  \item te: Sure. (Bubble increases, characters stay in the bubble, until the whole screen is in the bubble. Characters also move to the left lower corner.)
\end{enumerate}

\subsection{Define elliptic curves}
 \begin{enumerate}
   \item te: It deals with equations like $y^2=x^3-4x+1$, called elliptic curves (centred).
   \item st: What do they look like?
   \item te: Here (Equation moves up left, picture appears.)
   \item st: What are other examples?
   \item te: We need $y^2=x^3+$ integer $x+$ integers. (Equation and picture show other curves with coefficients framed)
   \item st: They are all symmetric ? (Back to $E$)
   \item te: Yes as $y$ only appears as an even power ($(-y)^2 = y^2 = \dots$)
 \end{enumerate}

 \subsection{Rational points}
 \begin{enumerate}
   \item st: What you study about them?
   \item te: I count solutions to the equation with both $x$ and $y$ rational. ($x,y\in\QQ$ appears)
   \item st: Hmm, if $x=0$ and $y$ \dots $y=1$ or $y=-1$.
   \item te: Yep. I chose an example with a lot of solutions (Points are highlighted successively, including one really complicated example)
   \item te: Some other elliptic curves have very few.
 \end{enumerate}

 \subsection{Projective curve}
 \begin{enumerate}
   \item te: We can write $x=X/Z$ and $y=Y/Z$ with integers where $Z$ is the least positive common denominator.
   \item (The equation transforms to substitute, move power).
   \item te: Multiply by $Z^3$. (Get to the projective equation $Y^2Z= X^3-4XZ^2+Z^3$)
   \item st: I know this is a projective curve.
   \item te: Indeed. We can now study integer solution $X$, $Y$, $Z$ to this equation instead.
   \item te: But if we, do we get extra solution, like when $Z$ is negative.
   We need to identify solutions when they are just scalar multiples of each other. ($(X,Y,Z) = (-X,-Y,-Z) = (2X,2Y,2Z)$ as $x = \frac{X}{Z} = {-X}{-Z} = {2X}{2Y}$)
   \item te: And we have added accidentality a new point, the one with $Z=0$ and $X=0$ and $Y=1$.
   To visual it as a point infinitely far away, we should move to see this. (Turn to the 3d Scene)
 \end{enumerate}

 % up to here



\section{Count rational points by height}
\begin{itemize}
  \item Define $H$ as max $X$ etc.
  \item Plot them logarithmically as towers on the points on a rank 2 curve
  \item Define $N(T)$
  \item Naive search finds them in that order
\end{itemize}

\section{Count modulo}
\begin{itemize}
  \item For an integer $U$, explain counting modulo $U$
  \item More rational points more points mod $U$ (is wrong)
  \item Set $U=T!$, really biiiiiig
\end{itemize}

\section{Conjecture and non-evidence}
\begin{itemize}
  \item Introduce quotient and make conjecture
  \item A convincing picture for rank 0 and one for rk 2
  \item A picture with a jump for a point
\end{itemize}

\section{Rank and original formulation}
\begin{itemize}
  \item Explain $N(T) \sim \log(T)^{r/2}$
  \item Say $r$ is rank of $E(\mathbb{Q})$
  \item Birch and Swinnerton-Dyer
\end{itemize}

\section{L-functions and GRH}
\begin{itemize}
  \item Say there is an $L$-function, Weil?
  \item Analytic continuation implied by our conjecture
  \item As well as GRH
  \item Write out limit?
\end{itemize}

\end{document}
